% !TeX root = ../../Book3.tex
\section{上半平面与模群预告}
\label{b3:sec:1104}

圆盘与上半平面是同一个共形模型。上半平面的自同构用实矩阵表示，既方便分类甲板变换，也给模群与模形式提供几何入口。

\subsection{上半平面}
\label{b3:subsec:110401}

上半平面记为 \(\mathbb H=\{z:\operatorname{Im}z>0\}\)。Cayley 变换 \(C(z)=(z-\mathrm{i})/(z+\mathrm{i})\) 将它双全纯映到圆盘，逆为 \(\mathrm{i}(1+w)/(1-w)\)。实轴加无穷远对应单位圆周，内部对称可在两种模型间共轭传递。

\ProseParagraphBreak
无穷远是球面边界点，不是半平面内部点。平移 \(z\mapsto z+1\) 没有内部固定点，却在边界固定无穷远；缩放 \(z\mapsto az\)、\(a>0\)、\(a\ne1\) 固定边界的0与无穷远。两种无固定点运动的商具有不同端部几何。

\subsection{Möbius 变换}
\label{b3:subsec:110402}

实系数分式 \(T(z)=(az+b)/(cz+d)\)、\(ad-bc>0\)，满足
\[
 \operatorname{Im}T(z)=\frac{(ad-bc)\operatorname{Im}z}{|cz+d|^2}>0.
\]
所以它是半平面自同构；反向由圆盘自同构分类与 Cayley 共轭得到所有自同构都如此。将矩阵整体除以行列式平方根，可令行列式为1；矩阵与其负矩阵给相同分式。

\ProseParagraphBreak
按实扩充轴上的固定点数，非恒等自同构分为椭圆型、抛物型、双曲型：分别有内部固定点、一个重边界固定点、两个不同边界固定点。对行列式1矩阵，判别式为 \((a+d)^2-4\)，故迹绝对值分别小于、等于、大于2。甲板变换自由，因而不能是非平凡椭圆型。

\subsection{\texorpdfstring{\(\mathrm{PSL}_2(\mathbb R)\)}{PSL2(R)}}
\label{b3:subsec:110403}

群 \(\operatorname{PSL}_2(\mathbb R)=\operatorname{SL}_2(\mathbb R)/\{\pm I\}\) 正是保定向半平面全纯自同构群。后面的双曲度量元 \(|\mathrm{d}z|/\operatorname{Im}z\) 在它作用下不变，可由上式与 \(T'(z)=1/(cz+d)^2\) 立即验证。

\ProseParagraphBreak
群的离散性是关于变换参数的拓扑，适当不连续性是关于在半平面的作用；对这里的离散子群，作用的局部性质可由矩阵紧性与内部点控制证明。要形成普通无分歧覆盖，还必须无挠，排除椭圆固定点。只写一个 \(\operatorname{PSL}_2(\mathbb R)\) 子群，并不足以自动得到光滑商曲面。

\ProseParagraphBreak
实 Möbius 变换的不动点类型可由二次方程核验。把矩阵归一化为行列式1，不动点方程的判别式为 \((a+d)^2-4\)。当其负时，两根共轭且一根在上半平面，是椭圆型；当其为0且不是恒等变换，只有一个实边界不动点，是抛物型；当其正时，有两个不同实边界不动点，是双曲型。\(\operatorname{tr}^2\) 而非迹本身才与 \(\mathrm{PSL}_2\) 中矩阵正负代表无关。

\ProseParagraphBreak
标准形分别是圆盘内旋转、上半平面平移和正伸缩。前者有内部不动点，后两者没有。覆盖的甲板群须自由，故不能含非平凡椭圆元。\(\mathrm{PSL}_2(\mathbb Z)\) 包含椭圆元，并不满足这个自由条件；完整模群商有特殊稳定子点，不能直接称为通常无分歧曲面覆盖。若转向没有稳定子的子群，才能按前面覆盖理论形成无分歧商；群的离散性和作用的自由性是不同条件。

\subsection{\texorpdfstring{\(\mathrm{PSL}_2(\mathbb Z)\)}{PSL2(Z)}}
\label{b3:subsec:110404}

模群 \(\operatorname{PSL}_2(\mathbb Z)\) 由整数矩阵给出，经典生成元为 \(T(z)=z+1\)、\(S(z)=-1/z\)。欧氏算法消去整数矩阵的下左元素，可将任意行列式1整数矩阵写成这两类操作的乘积，故它们生成模群。

\ProseParagraphBreak
\(S\) 固定 \(\mathrm{i}\)，\(ST\) 固定 \((-1+\sqrt3\,\mathrm{i})/2\)。因此整个模群商具有椭圆点，商投影不是处处无分歧覆盖。通常基本区域 \(|\operatorname{Re}z|\le1/2,|z|\ge1\) 的两竖边以 \(T\) 识别，圆弧边以 \(S\) 识别；无穷远对应尖点。模形式将要求全纯函数在变换下按 \((cz+d)^k\) 的权变换，并另检验尖点展开。这里的几何预告不先假定模形式的维数或傅里叶展开定理。
